\lyseisE{Επανάληψη κεφαλαίου Β3}
\solr{Β3}{1} $25^\circ$ οξεία, $90^\circ$ ορθή, $135^\circ$ αμβλεία, $180^\circ$ ευθεία, $300^\circ$ μη κυρτή, $360^\circ$ πλήρης.
\solr{Β3}{2} $\mathrm{ΜΒ}=7$ cm, γιατί κάθε σημείο της μεσοκαθέτου ισαπέχει από τα άκρα. Σχεδιάζουμε δύο κύκλους ίδιας ακτίνας (μεγαλύτερης από το μισό του $\mathrm{ΑΒ}$) με κέντρα $\mathrm{Α}$, $\mathrm{Β}$, και φέρνουμε την ευθεία των δύο σημείων τομής τους.
\solr{Β3}{3} $x\hat{\mathrm{Ο}}\delta=112^\circ:2=56^\circ$. Τα σημεία της διχοτόμου ισαπέχουν από τις πλευρές, άρα απέχει $3$ cm και από την $\mathrm{Ο}y$.
\solr{Β3}{4} Ισοσκελές (δύο ίσες πλευρές). Περίμετρος $9+9+12=30$ cm.
\solr{Β3}{5} $\hat{\Gamma}=\hat{\mathrm{Β}}=65^\circ$. Η διάμεσος προς τη βάση: $\mathrm{ΒΜ}=10:2=5$ cm, και είναι και ύψος: $\mathrm{ΑΜ}\hat{\mathrm{Β}}=90^\circ$.
\solr{Β3}{6} $64^\circ$ (κατακορυφήν), και δύο γωνίες $180^\circ-64^\circ=116^\circ$.
\solr{Β3}{7} Συμπληρωματική: $90^\circ-38^\circ=52^\circ$. Παραπληρωματική: $180^\circ-38^\circ=142^\circ$.
\solr{Β3}{8} Είναι παράλληλες (κάθετες στην ίδια ευθεία). Από σημείο εκτός της $\varepsilon_1$ περνά μία μόνο παράλληλη προς αυτήν.
\solr{Β3}{9} Οι οξείες γωνίες είναι όλες $70^\circ$ και οι αμβλείες όλες $180^\circ-70^\circ=110^\circ$ (τέσσερις γωνίες η καθεμία).
\solr{Β3}{10} $180^\circ-(58^\circ+64^\circ)=58^\circ$. Οξυγώνιο (και ισοσκελές).
\solr{Β3}{11} $(180^\circ-36^\circ):2=72^\circ$ η καθεμία.
\solr{Β3}{12} $x+(x+20^\circ)=90^\circ$, $2x=70^\circ$, $x=35^\circ$. Οι γωνίες είναι $35^\circ$ και $55^\circ$.
\solr{Β3}{13} Περίμετρος $2\cdot(9+6)=30$ cm. Οι γωνίες είναι $75^\circ$, $105^\circ$, $75^\circ$, $105^\circ$.
\solr{Β3}{14} Οι γωνίες της μικρής βάσης είναι $180^\circ-58^\circ=122^\circ$ η καθεμία. Άρα $58^\circ$, $58^\circ$, $122^\circ$, $122^\circ$.
\solr{Β3}{15} Περίμετρος $4\cdot7{,}5=30$ cm. Οι διαγώνιοι διχοτομούνται και τέμνονται κάθετα.
\solr{Β3}{16} $\mathrm{ΒΔ}=15$ cm (ίσες διαγώνιοι). Οι διαγώνιοι διχοτομούνται, άρα η απόσταση είναι $15:2=7{,}5$ cm.
\solr{Β3}{17} Διάμετρος $2\cdot4=8$ cm. Το $\mathrm{Α}$ απέχει $5>4$ από το κέντρο, άρα είναι έξω από τον κύκλο.
\solr{Β3}{18} Για $3$ cm: δύο κοινά σημεία (τέμνουσα). Για $5$ cm: ένα (εφαπτομένη). Για $7$ cm: κανένα (εξωτερική).
\solr{Β3}{19} α) Λάθος. β) Σωστό. γ) Σωστό. δ) Σωστό. ε) Λάθος: δύο ύψη βρίσκονται εκτός.
\solr{Β3}{20} α) Σωστό. β) Λάθος: οι διαγώνιοι του ρόμβου είναι κάθετες, όχι πάντα ίσες. γ) Σωστό. δ) Λάθος: δύο μόνο. ε) Σωστό.
